\clearpage
\thispagestyle{empty}
\pdfbookmark[0]{Allgemeine Theorie des Erdmagnetismus: Ausgabe und Quellen}{erdmagnetismus-source-notice}
\begin{center}\Large\textbf{Allgemeine Theorie des Erdmagnetismus}\end{center}

KI-Übertragung des deutschen Textes aus Carl Friedrich Gauss, \textit{Werke}, Band V (Göttingen, 1867), S.~119--193, in modernes Deutsch mit KI-erstelltem Satz. Die Ausgabe enthält beide Titel, die Einleitung, alle 41 Artikel und acht historische Anmerkungen, den Nachtrag mit seinen historischen redaktionellen Passagen sowie die Vergleichs- und Hilfstafeln. Die beiden Diagramme der historischen Tafel sind als bearbeitbare Nachzeichnungen mit angenäherter Schraffur und Linienstärke wiedergegeben.

Die im Text genannte historische Karte für die Werte von \(V/R\) gehört nicht zur verfügbaren bearbeitbaren Vorlage. Die drei Abteilungen der in Artikel 30 beschriebenen Karte sind in der Ausgabe von Glassmeier und Tsurutani nach den Tafeln des Ergänzungsbandes von 1839 abgebildet.

Die englische Übersetzung von Elizabeth Juliana Sabine, revidiert von Sir John Herschel, erschien in Richard Taylors \textit{Scientific Memoirs}, Band 2 (1841), S.~184--251. Sie und die revidierte Übersetzung von K.-H.~Glassmeier und B.~T.~Tsurutani in \textit{History of Geo- and Space Sciences} 5 (2014), S.~11--62, wurden zum Vergleich herangezogen. Für die heutige deutsche Terminologie dienten die IGRF-Erläuterungen des GFZ als Vergleich.

\noindent\href{https://gdz.sub.uni-goettingen.de/id/PPN236006339}{Deutsche Vorlage (1867)}\quad
\href{https://en.wikisource.org/wiki/Scientific_Memoirs/2/General_Theory_of_Terrestrial_Magnetism?oldid=4352379}{Sabine und Herschel (1841)}\quad
\href{https://doi.org/10.5194/hgss-5-11-2014}{Glassmeier und Tsurutani (2014)}\quad
\href{https://isdc.gfz.de/igrf-declination-calculator/erlaeuterungen-de/}{GFZ}
\clearpage
